% GENERATED FILE. Edit canonical lessons, then run npm run generate:books.
% canonical-source-hash: fnv1a64:10d978b0130efbab
% canonical-lessons: HI-C513-silna, HI-C513-silvana, HI-C513-katarna, HI-C513-pirona, HI-C513-udharna, HI-R513-review-506-513

\chapter{At the Tailor's}
\label{ch:hi-at-the-tailor-s}
\hlchaptermodality{513}

\begin{chapteropening}
\textbf{By the end of this chapter:} \emph{I can ask a tailor to stitch, trim and unpick clothes.}

Complete the last lesson of chapter 513: I can ask a tailor to stitch, trim and unpick clothes.
\end{chapteropening}

\section[\texorpdfstring{\hi{सिलना} (silnā)}{सिलना (silnā)}]{\hi{सिलना} (silnā) — to be stitched; to get sewn}

\label{lesson:HI-C513-silna}



\begin{quote}
Before the new one: say the Hindi for a loom, then the Hindi for a knot.
\end{quote}

\subsection*{You'll want to know: \hi{सिलना}}
\textbf{\hi{सिलना}} — \emph{silnā} — \textquotedblleft{}to be stitched; to get sewn\textquotedblright{}.

\begin{quote}
\hi{मेरी} \hi{नई} \hi{कमीज़} \hi{कल} \hi{तक} \hi{सिल} \hi{जाएगी।}
\end{quote}

\begin{quote}
\emph{Merī naī kamīz kal tak sil jāegī.}
\end{quote}

\begin{quote}
My new shirt will be stitched by tomorrow.
\end{quote}

Past (it …): \hi{वह} \hi{कुर्ता} \hi{सिल} \hi{गया}, \emph{vah kurtā sil gayā}.

\begin{cousinweb}
\textbf{\hi{सिलना}} (\emph{silnā}) is the intransitive partner of \emph{sīnā}, \textquotedblleft{}to sew\textquotedblright{}: CDIAL files both under Sanskrit \emph{sīvyati}, \textquotedblleft{}sews\textquotedblright{} (CDIAL 13444).
\end{cousinweb}

\begin{grammarlens}[title={What you've built}]
One more word for describing what you lived and felt.
\end{grammarlens}

\subsection*{Your turn}
\begin{itemize}
\raggedright
  \item \emph{Say it:} \emph{silnā}
  \item \emph{Say it:} \emph{silnā}, once more
  \item \emph{Say it:} \emph{Merī naī kamīz kal tak sil jāegī.} — My new shirt will be stitched by tomorrow.
  \item \emph{Recall:} say the Hindi for a loom, then the Hindi for a knot, then say \emph{silnā} again
\end{itemize}

\begin{tcolorbox}[breakable,colback=teal!4,colframe=teal!35!black,title={Before you move on}]
What does \hi{सिलना} mean? (\textquotedblleft{}To be stitched; to get sewn\textquotedblright{}.) Say it once more.
\end{tcolorbox}
\section[\texorpdfstring{\hi{सिलवाना} (silvānā)}{सिलवाना (silvānā)}]{\hi{सिलवाना} (silvānā) — to have (clothes) stitched}

\label{lesson:HI-C513-silvana}



\begin{quote}
Before the new one: say the Hindi for a knot, then the Hindi for to be stitched.
\end{quote}

\subsection*{You'll want to know: \hi{सिलवाना}}
\textbf{\hi{सिलवाना}} — \emph{silvānā} — \textquotedblleft{}to have (clothes) stitched\textquotedblright{}.

\begin{quote}
\hi{मैंने} \hi{दर्ज़ी} \hi{से} \hi{एक} \hi{नया} \hi{कुर्ता} \hi{सिलवाया।}
\end{quote}

\begin{quote}
\emph{Maiṁne darzī se ek nayā kurtā silvāyā.}
\end{quote}

\begin{quote}
I had a new kurta stitched by the tailor.
\end{quote}

Past, with \hi{ने} (he or she …): \hi{उसने} \hi{एक} \hi{कमीज़} \hi{सिलवाई}, \emph{usne ek kamīz silvāī}.

\begin{cousinweb}
\textbf{\hi{सिलवाना}} (\emph{silvānā}) is the second causative of \emph{sīnā}, \textquotedblleft{}to sew\textquotedblright{}: CDIAL files \emph{sīnā} and its causative \emph{silānā} under Sanskrit \emph{sīvyati}, \textquotedblleft{}sews\textquotedblright{} (CDIAL 13444). It means getting someone else to do the sewing.
\end{cousinweb}

\begin{grammarlens}[title={What you've built}]
One more word for describing what you lived and felt.
\end{grammarlens}

\subsection*{Your turn}
\begin{itemize}
\raggedright
  \item \emph{Say it:} \emph{silvānā}
  \item \emph{Say it:} \emph{silvānā}, once more
  \item \emph{Say it:} \emph{Maiṁne darzī se ek nayā kurtā silvāyā.} — I had a new kurta stitched by the tailor.
  \item \emph{Recall:} say the Hindi for a knot, then the Hindi for to be stitched, then say \emph{silvānā} again
\end{itemize}

\begin{tcolorbox}[breakable,colback=teal!4,colframe=teal!35!black,title={Before you move on}]
What does \hi{सिलवाना} mean? (\textquotedblleft{}To have (clothes) stitched\textquotedblright{}.) Say it once more.
\end{tcolorbox}
\section[\texorpdfstring{\hi{कतरना} (katarnā)}{कतरना (katarnā)}]{\hi{कतरना} (katarnā) — to snip; to trim (with scissors)}

\label{lesson:HI-C513-katarna}



\begin{quote}
Before the new one: say the Hindi for to be stitched, then the Hindi for to have (clothes) stitched.
\end{quote}

\subsection*{You'll want to know: \hi{कतरना}}
\textbf{\hi{कतरना}} — \emph{katarnā} — \textquotedblleft{}to snip; to trim (with scissors)\textquotedblright{}.

\begin{quote}
\hi{उसने} \hi{कैंची} \hi{से} \hi{कपड़े} \hi{के} \hi{कोने} \hi{कतरे।}
\end{quote}

\begin{quote}
\emph{Usne kaĩcī se kapṛe ke kone katre.}
\end{quote}

\begin{quote}
She snipped the corners of the cloth with scissors.
\end{quote}

Past, with \hi{ने} (he or she …): \hi{उसने} \hi{धागा} \hi{कतरा}, \emph{usne dhāgā katrā}.

\begin{cousinweb}
\textbf{\hi{कतरना}} (\emph{katarnā}) is a tadbhava, inherited from Sanskrit \emph{kartarī}, \textquotedblleft{}scissors\textquotedblright{} (CDIAL 2858).
\end{cousinweb}

\begin{grammarlens}[title={What you've built}]
One more word for describing what you lived and felt.
\end{grammarlens}

\subsection*{Your turn}
\begin{itemize}
\raggedright
  \item \emph{Say it:} \emph{katarnā}
  \item \emph{Say it:} \emph{katarnā}, once more
  \item \emph{Say it:} \emph{Usne kaĩcī se kapṛe ke kone katre.} — She snipped the corners of the cloth with scissors.
  \item \emph{Recall:} say the Hindi for to be stitched, then the Hindi for to have (clothes) stitched, then say \emph{katarnā} again
\end{itemize}

\begin{tcolorbox}[breakable,colback=teal!4,colframe=teal!35!black,title={Before you move on}]
What does \hi{कतरना} mean? (\textquotedblleft{}To snip; to trim (with scissors)\textquotedblright{}.) Say it once more.
\end{tcolorbox}
\section[\texorpdfstring{\hi{पिरोना} (pironā)}{पिरोना (pironā)}]{\hi{पिरोना} (pironā) — to thread (a needle); to string (beads)}

\label{lesson:HI-C513-pirona}



\begin{quote}
Before the new one: say the Hindi for to have (clothes) stitched, then the Hindi for to snip.
\end{quote}

\subsection*{You'll want to know: \hi{पिरोना}}
\textbf{\hi{पिरोना}} — \emph{pironā} — \textquotedblleft{}to thread (a needle); to string (beads)\textquotedblright{}.

\begin{quote}
\hi{दर्ज़ी} \hi{ने} \hi{सूई} \hi{में} \hi{धागा} \hi{पिरोया।}
\end{quote}

\begin{quote}
\emph{Darzī ne sūī meṁ dhāgā piroyā.}
\end{quote}

\begin{quote}
The tailor threaded the needle.
\end{quote}

Past, with \hi{ने} (he or she …): \hi{उसने} \hi{माला} \hi{में} \hi{फूल} \hi{पिरोए}, \emph{usne mālā meṁ phūl piroe}.

\begin{cousinweb}
\textbf{\hi{पिरोना}} (\emph{pironā}) is a tadbhava, inherited from Sanskrit \emph{parivayati}, \textquotedblleft{}weaves round\textquotedblright{} (CDIAL 7869). CDIAL gives the Hindi verb as \textquotedblleft{}to plait, string (pearls), thread a needle\textquotedblright{}.
\end{cousinweb}

\begin{grammarlens}[title={What you've built}]
One more word for describing what you lived and felt.
\end{grammarlens}

\subsection*{Your turn}
\begin{itemize}
\raggedright
  \item \emph{Say it:} \emph{pironā}
  \item \emph{Say it:} \emph{pironā}, once more
  \item \emph{Say it:} \emph{Darzī ne sūī meṁ dhāgā piroyā.} — The tailor threaded the needle.
  \item \emph{Recall:} say the Hindi for to have (clothes) stitched, then the Hindi for to snip, then say \emph{pironā} again
\end{itemize}

\begin{tcolorbox}[breakable,colback=teal!4,colframe=teal!35!black,title={Before you move on}]
What does \hi{पिरोना} mean? (\textquotedblleft{}To thread (a needle); to string (beads)\textquotedblright{}.) Say it once more.
\end{tcolorbox}
\section[\texorpdfstring{\hi{उधड़ना} (udhaṛnā)}{उधड़ना (udhaṛnā)}]{\hi{उधड़ना} (udhaṛnā) — to come unstitched; to come apart}

\label{lesson:HI-C513-udharna}



\begin{quote}
Before the new one: say the Hindi for to snip, then the Hindi for to thread (a needle).
\end{quote}

\subsection*{You'll want to know: \hi{उधड़ना}}
\textbf{\hi{उधड़ना}} — \emph{udhaṛnā} — \textquotedblleft{}to come unstitched; to come apart\textquotedblright{}.

\begin{quote}
\hi{मेरी} \hi{जेब} \hi{उधड़} \hi{गई।}
\end{quote}

\begin{quote}
\emph{Merī jeb udhaṛ gaī.}
\end{quote}

\begin{quote}
My pocket has come unstitched.
\end{quote}

Past (it …): \hi{वह} \hi{कमीज़} \hi{उधड़} \hi{गई}, \emph{vah kamīz udhaṛ gaī}.

\begin{cousinweb}
\textbf{\hi{उधड़ना}} (\emph{udhaṛnā}) is a tadbhava, inherited from Sanskrit \emph{uddharati}, \textquotedblleft{}takes out\textquotedblright{} (CDIAL 2009). CDIAL gives both Hindi verbs as \textquotedblleft{}to be taken out, be unravelled\textquotedblright{}; \emph{udhaṛnā} is the one that happens by itself.
\end{cousinweb}

\begin{grammarlens}[title={What you've built}]
One more word for describing what you lived and felt.
\end{grammarlens}

\subsection*{Your turn}
\begin{itemize}
\raggedright
  \item \emph{Say it:} \emph{udhaṛnā}
  \item \emph{Say it:} \emph{udhaṛnā}, once more
  \item \emph{Say it:} \emph{Merī jeb udhaṛ gaī.} — My pocket has come unstitched.
  \item \emph{Recall:} say the Hindi for to snip, then the Hindi for to thread (a needle), then say \emph{udhaṛnā} again
\end{itemize}

\begin{tcolorbox}[breakable,colback=teal!4,colframe=teal!35!black,title={Before you move on}]
What does \hi{उधड़ना} mean? (\textquotedblleft{}To come unstitched; to come apart\textquotedblright{}.) Say it once more.
\end{tcolorbox}
\section[(dialogue)]{Looking back: from The Craftsman's Bench to At the Tailor's}

\label{lesson:HI-R513-review-506-513}



\begin{quote}
Say the Hindi for a craftsman and a tool.
\end{quote}

\subsection*{Your turn}
Say these aloud:

\begin{itemize}
\raggedright
  \item a craftsman, a tool, a skill, to shape, a mechanic
  \item a colleague, an officer, a master (of a craft or an art), skill, talent
  \item an interview (for a job or a place), an appointment (to a post), an application (formal), admission (to a school or course), a salary, pay (formal)
  \item a strike (stopping work), a resignation, a contribution, high-handedness, a slogan
  \item a carpenter's plane, a chisel, to rub, to wear down, to knock in, to bore
  \item a potter's wheel, to mix (with water), a potter's kiln, a mould, to pour into a mould
  \item cotton (the plant, the raw fibre), a spinning wheel, to spin (thread), a loom, a knot
  \item to be stitched, to have (clothes) stitched, to snip, to thread (a needle), to come unstitched
\end{itemize}

\begin{tcolorbox}[breakable,colback=teal!4,colframe=teal!35!black,title={Before you move on}]
A craftsman? (\textbf{\hi{कारीगर}}, \emph{kārīgar}.) A tool? (\textbf{\hi{औज़ार}}, \emph{auzār}.) A skill? (\textbf{\hi{हुनर}}, \emph{hunar}.) To shape? (\textbf{\hi{गढ़ना}}, \emph{gaṛhnā}.) A mechanic? (\textbf{\hi{मिस्त्री}}, \emph{mistrī}.) A colleague? (\textbf{\hi{सहकर्मी}}, \emph{sahkarmī}.) An officer? (\textbf{\hi{अफ़सर}}, \emph{afsar}.) A master (of a craft or an art)? (\textbf{\hi{उस्ताद}}, \emph{ustād}.) Skill? (\textbf{\hi{कौशल}}, \emph{kauśal}.) Talent? (\textbf{\hi{प्रतिभा}}, \emph{pratibhā}.) An interview (for a job or a place)? (\textbf{\hi{इंटरव्यू}}, \emph{iṇṭarvyū}.) An appointment (to a post)? (\textbf{\hi{नियुक्ति}}, \emph{niyukti}.) An application (formal)? (\textbf{\hi{आवेदन}}, \emph{āvedan}.) Admission (to a school or course)? (\textbf{\hi{दाख़िला}}, \emph{dākhilā}.) A salary, pay (formal)? (\textbf{\hi{वेतन}}, \emph{vetan}.) A strike (stopping work)? (\textbf{\hi{हड़ताल}}, \emph{haṛtāl}.) A resignation? (\textbf{\hi{इस्तीफ़ा}}, \emph{istīfā}.) A contribution? (\textbf{\hi{चंदा}}, \emph{candā}.) High-handedness? (\textbf{\hi{मनमानी}}, \emph{manmānī}.) A slogan? (\textbf{\hi{नारा}}, \emph{nārā}.) A carpenter's plane? (\textbf{\hi{रंदा}}, \emph{randā}.) A chisel? (\textbf{\hi{छेनी}}, \emph{chenī}.) To rub, to wear down? (\textbf{\hi{घिसना}}, \emph{ghisnā}.) To knock in? (\textbf{\hi{ठोकना}}, \emph{ṭhoknā}.) To bore? (\textbf{\hi{छेदना}}, \emph{chednā}.) A potter's wheel? (\textbf{\hi{चाक}}, \emph{cāk}.) To mix (with water)? (\textbf{\hi{सानना}}, \emph{sānnā}.) A potter's kiln? (\textbf{\hi{आँवा}}, \emph{ā̃vā}.) A mould? (\textbf{\hi{साँचा}}, \emph{sā̃cā}.) To pour into a mould? (\textbf{\hi{ढालना}}, \emph{ḍhālnā}.) Cotton (the plant, the raw fibre)? (\textbf{\hi{कपास}}, \emph{kapās}.) A spinning wheel? (\textbf{\hi{चरखा}}, \emph{carkhā}.) To spin (thread)? (\textbf{\hi{कातना}}, \emph{kātnā}.) A loom? (\textbf{\hi{करघा}}, \emph{karghā}.) A knot? (\textbf{\hi{गाँठ}}, \emph{gā̃ṭh}.) To be stitched? (\textbf{\hi{सिलना}}, \emph{silnā}.) To have (clothes) stitched? (\textbf{\hi{सिलवाना}}, \emph{silvānā}.) To snip? (\textbf{\hi{कतरना}}, \emph{katarnā}.) To thread (a needle)? (\textbf{\hi{पिरोना}}, \emph{pironā}.) To come unstitched? (\textbf{\hi{उधड़ना}}, \emph{udhaṛnā}.)
\end{tcolorbox}
